\subsection{Structural Unification of the Kinematic Deficit}
  \label{subsec:structural-unification-of-the-kinematic-deficit}

  The transition from galactic scales to cosmological expansion is not a change in
  the gravitational law but a property of the effective inference metric itself.

  Equation~(1) expresses a bounded-response relation between the Newtonian source
  $g_N$ and the inferred effective acceleration $g_{\mathrm{eff}}$. In the deep
  saturation regime, the response function $\mu(x)$ enforces

  \[
    g_{\mathrm{eff}} \propto \sqrt{a_\star g_N},
  \]

  which limits the admissible growth of the inferred acceleration when
  $g_N \ll a_\star$.

  This limitation is not tied to radial geometry.
  It reflects a bound on the norm of the inferred gradient of the effective potential,
  \[
    \|\nabla \Phi_{\mathrm{eff}}\| \leq a_\star,
  \]
  which applies independently of whether the dilution is radial
  (galactic outskirts) or volumetric (cosmic voids).

  In galactic systems, decreasing surface density drives $g_N$ below $a_\star$,
  producing flat rotation curves as shown in Section~3.
  In cosmological voids, decreasing volumetric density suppresses the effective
  gradient of the large-scale metric inference, leading to a bounded deficit in
  the locally inferred expansion rate.

  The projection deficit $\delta$ entering Eq.~(8) is therefore the volumetric
  analogue of the radial deficit encoded in Eq.~(4).
  Both emerge from the same bounded-response constraint on the effective metric.

  The relationship between Eq.~(1) and Eq.~(8) is thus mediated by a scale-invariant
  inference threshold $a_\star$, which remains constant across geometries.
